\documentclass[12pt,aspectratio=169]{beamer}
\input{header_beamer}
\usetheme{Boadilla}
\usecolortheme{beaver}
\setbeamertemplate{page number in head/foot}{}
\setbeamertemplate{navigation symbols}{}

\title{Lecture-05: Random Vectors}
\date{}

\begin{document}
\pagestyle{empty}
\maketitle
\begin{frame}{Random vectors} 
\begin{Definition}[Projection]<1->
For a vector $x \in \R^n$, we can define $\pi_i: \R^n \to \R$ is the \textbf{projection} of an $n$-length vector onto its $i$-th component, such that $\pi_i(x) = x_i$. 
\end{Definition}
\begin{Definition}[Borel sigma algebra]<2->
\begin{itemize}
\item The Borel sigma algebra over space $\R^n$ is the smallest sigma algebra generated by the family $(\pi_i^{-1}(B_x): x \in \R, i \in [n])$ denoted by $\sB(\R^n)$. 
\item The elements of the Borel sigma algebra are called Borel sets. 
\end{itemize}
\end{Definition}
\end{frame}
\begin{frame}{Random vectors}
\begin{block}{Reminder}
\begin{itemize}
\item <1-> Projection $\pi_i: \R^n \to \R$ is a Borel measurable function for all $i \in [n]$. 
\item <2-> For a subset $A \subseteq \R$ and projection $\pi_i: \R^n \to \R$, 
\EQ{
\pi_i^{-1}(A) = \set{x \in \R^n: x_i \in A} = \R \times \dots \times A \times \dots \times \R.
}
\item <3-> Thus, for any $A \in \sB(\R)$, we have $\pi_i^{-1}(A) \in \sB(\R^n)$.
\end{itemize}
\end{block}
\end{frame}
\begin{frame}{Random vectors}
\begin{Definition}[Random vectors] 
\begin{itemize}
\item <1-> Consider a probability space $(\Omega, \sF, P)$ and a finite $n \in \N$.\\ 
\item <2-> A \textbf{random vector} $X: \Omega \to \R^n$ is an $\sF$-measurable mapping from the sample space to an $n$-length real-valued vector.\\
\item <3-> That is, for any $x \in \R^n$, we have %the event 
\EQ{
A_X(x) \triangleq \set{\omega \in \Omega: X_1(\omega) \le x_1, \dots, X_n(\omega) \le x_n} 
= \cap_{i=1}^nX_i^{-1}(-\infty, x_i] 
\in \sF.
}
\end{itemize}
\end{Definition}
\end{frame}
\begin{frame}{Random vectors}
\begin{Example}[Tuple of indicators] 
Consider $(\Omega, \sF, P)$, a finite $n \in \N$, and events $A_1, \dots, A_n \in \sF$. \\
Define $X: \Omega \to \set{0,1}^n$ by $X_i(\omega) \triangleq \Ind{A_i}(\omega)$ for all outcomes $\omega \in \Omega$.\\
Let $x \in \R^n$, then $A_X(x) = \cap_{i=1}^n\Ind{A_i}^{-1}(-\infty, x_i]$. 
\EQ{
\Ind{A_i}^{-1}(-\infty, x_i] = 
\begin{cases}
\Omega, & x_i \ge 1,\\
A_i^c, & x_i \in [0,1),\\
\emptyset, &x_i < 0.
\end{cases}
}
It follows that $A_X(x)$ lies in $\sF$, and hence $X$ is an $\sF$-measurable random vector. 
\end{Example}
\end{frame}

\begin{frame}{Random vectors}
\begin{block}{Theorem}
Consider a probability space $(\Omega, \sF, P)$, and a finite $n \in \N$. 
A mapping $X: \Omega \to \R^n$ is a random vector if and only if $X_i \triangleq \pi_i \circ X: \Omega \to \R$ are random variables for all $i \in [n]$. 
\end{block}
\end{frame}
\begin{frame}{Random vectors}
\begin{Proof}
\begin{itemize}
\item <1-> For any $i \in [n]$ and $x_i \in \R$, we take $x = (\infty, \dots, x_i, \dots, \infty)$.\\
\implies 
$\pi_i^{-1}(-\infty, x_i] = \R \times \dots \times (-\infty, x_i] \times \dots \times \R \in \sB(\R^n).$\\
Define $A_{X_i}(x_i) \triangleq X_i^{-1}(-\infty, x_i]$, then
\EQN{
\label{eqn:RandomVector:MarginalEvent}
A_X(x) = \cap_{j=1}^nX_j^{-1}(-\infty, x_j]
= X_i^{-1}(-\infty, x_i] 
= A_{X_i}(x_i) \in \sF.  
} 

\item <2-> For any $x \in \R^n$, we have $A_{X_i}(x_i) = X_i^{-1}(-\infty,x_i] \in \sF$ for all $i \in [n]$\\
\EQN{
\label{eqn:RandomVector:JointEvent}
A_X(x) = \cap_{i=1}^nA_{X_i}(x_i) \in \sF.
} 
\end{itemize}
\end{Proof}
\end{frame}
\begin{frame}{Distribution of random vectors}
\begin{Definition}[Joint distribution] 
Consider $(\Omega, \sF, P)$ and a finite $n \in \N$. 
The \textbf{joint distribution function} of a random vector $X: \Omega \to \R^n$ is defined as the mapping $F_X: \R^n \to [0,1]$ such that 
\EQ{
F_X(x) 
\triangleq P(A_X(x)) 
= P(\cap_{i=1}^nA_{X_i}(x_i)).
}  
\end{Definition}
\end{frame}
\begin{frame}{Distribution of random vectors}
\begin{Example}[Tuple of indicators] 
\begin{itemize}
\item <1-> Consider $(\Omega, \sF, P)$, a finite $n \in \N$, and events $A_1, \dots, A_n \in \sF$, 
that define  $X \triangleq (\Ind{A_1}, \dots, \Ind{A_n})$.\\ 
\item <2-> For any $x \in \R^n$, define index sets $I_0(x) \triangleq \set{i \in [n]: x_i < 0}$ and $I_1(x) \triangleq \set{i \in [n]: x_i \in [0,1)}$  
\item <3-> \EQ{
F_X(x) = 
\begin{cases}
1, & I_0(x)\cup I_1(x) = \emptyset,\\
P(\cap_{i \in I_1(x)}A_i^c), & I_0(x) = \emptyset, I_1(x) \neq \emptyset,\\ 
0, & I_0(x) \neq \emptyset.
\end{cases}
}
\end{itemize}
\end{Example}
\end{frame}

\begin{frame}{Distribution of random vectors}
\begin{Definition}[Marginal distribution]<1->
For a random vector $X: \Omega \to \R^n$ defined on $(\Omega, \sF, P)$ and $i \in [n]$,  
the distribution of $X_i \triangleq \pi_i\circ X: \Omega \to \R$ is called the $i$th \textbf{marginal distribution}, and denoted by $F_{X_i}: \Omega \to [0,1]$. % for all $i \in [n]$. 
\end{Definition}


\begin{Example}[Tuple of indicators]<2->
Consider $(\Omega, \sF, P)$, a finite $n \in \N$, and events $A_1, \dots, A_n \in \sF$, 
that define $X \triangleq (\Ind{A_1}, \dots, \Ind{A_n})$.\\
The $i$th marginal distribution is given by  
\EQ{
F_{X_i}(x) = (1-P(A_i))\Ind{[0,1)}(x) + \Ind{[1, \infty)}(x).
}
\end{Example}
\end{frame}

\begin{frame}{Distribution of random vectors}
\begin{cor}[Marginal distribution]<1->
Consider $X: \Omega \to \R^n$ defined  $(\Omega, \sF, P)$ with $F_X: \R^n \to [0,1]$.\\ 
The $i$th \textbf{marginal distribution} and can be obtained from the joint distribution of $X$ as 
\EQ{
F_{X_i}(x_i) = \lim_{x_j \to \infty, \text{ for all }j \neq i}F_X(x). 
}
\end{cor}
\begin{Proof}<2->
For any $i \in [n]$ and $x_i \in \R$,
we have $X_i^{-1}(-\infty, x_i] = A_X(x)$ for $x = (\infty, \dots, x_i, \dots, \infty)$ 
from~\eqref{eqn:RandomVector:MarginalEvent}.  
\end{Proof}
\end{frame}
\begin{frame}{Properties of the joint distribution function}
\begin{block}{Lemma} 
Consider $X: \Omega \to \R^n$ defined on $(\Omega,\sF,P)$. 
$F_X: \R^n \to [0,1]$ satisfies the following properties. 
\begin{enumerate}[(i)]
\item <2-> For $x, y \in \R^n$ such that $x_i \le y_i$ for each $i \in [n]$, we have $F_X(x) \le F_X(y)$. 
\item <3-> The function $F_X(x)$ is right continuous at all points $x \in \R^n$. 
\item <4-> The lower limit is $\lim_{x_i \to -\infty}F_X(x) = 0$, and the upper limit is $\lim_{x_i \to \infty, i \in [n]}F_X(x) = 1$.
\end{enumerate}
\end{block}
\end{frame}
\begin{frame}{Properties of the joint distribution function}
\begin{Proof}
Consider $X: \Omega \to \R^n$ defined on $(\Omega, \sF, P)$ and any $x \in \R^n$. 
\begin{enumerate}[(i)]
\item <2-> We can verify that $A_X(x)  = \cap_{i=1}^nA_{X_i}(x_i) \subseteq \cap_{i=1}^nA_{X_i}(y_i) = A(y)$. 
\item <3-> The proof is similar to the proof for single random variable.  
\item <4-> The event $A_X(x) = \emptyset$ when $x_i = -\infty$ for some $i \in [n]$ and $A_X(x) = \Omega$ when $x_i = \infty$ for all $i \in [n]$, 
hence the result follow. 
\end{enumerate}
\end{Proof}
\end{frame}
\begin{frame}{Joint distribution function}
\begin{Example}[Probability of rectangular events] 
Consider $(\Omega, \sF, P)$ and a random vector $X: \Omega \to \R^2$. 
Consider the points $x \le y\in \R^2$ and the events 
\eq{
&B_1 \triangleq \set{x_1 < X_1 \le y_1} = A_{X_1}(y_1)\setminus A_{X_1}(x_1)\in \sF, \\
&B_2 \triangleq \set{x_2 < X_2 \le y_2} = A_{X_2}(y_2)\setminus A_{X_2}(x_2)\in \sF.
} 
The marginal probabilities are given by
\eq{
&P(B_1) 
= P(A_{X_1}(y_1))-P(A_{X_1}(x_1)) 
= F_{X_1}(y_1) - F_{X_1}(x_1),\\
&P(B_2) 
= P(A_{X_2}(y_2))-P(A_{X_2}(x_2)) 
= F_{X_2}(y_2) - F_{X_2}(x_2).
}
\end{Example}
\end{frame}

\begin{frame}{Example: Probability of rectangular events}
\begin{Example}[Continued]
Writing $x = (x_1,x_2)$ and $y = (y_1,y_2)$, 
the end points of the rectangular event $B_1\cap B_2$ are points $x, (y_1,x_2), y, (x_1, y_2)$.
\EQ{
B_1\cap B_2  = (A_X(y) \setminus A_X(x_1,y_2)) \setminus (A_X(y_1,x_2)\setminus A_X(x)) \in \sF.
}
Hence, we can write the probability of this rectangular event as 
\EQ{
P(B_1\cap B_2) = (F_X(y)-F_X(x_1,y_2)) - (F_X(y_1,x_2)-F_X(x)).
}
\end{Example}
\end{frame}

\begin{frame}{Event space generated by random vectors}
\begin{Definition} %[Event space generated by random vectors] 
Consider $(\Omega, \sF, P)$ and a finite $n \in \N$.\\
The \textbf{event space generated by a random vector $X:\Omega \to \R^n$} is the smallest $\sigma$-algebra generated by the collection of events $(A_X(x): x \in \R^n)$ and denoted by $\sigma(X) \triangleq \sigma(A_X(x):x \in \R^n)$.  
\end{Definition}
\end{frame}
\begin{frame}{Event space generated by random vectors}
\begin{block}{Theorem}
Consider a probability space $(\Omega, \sF, P)$, a finite $n \in \N$, 
a random vector $X: \Omega \to \R^n$, and its projections $X_i \triangleq \pi_i \circ X$ for all $i \in [n]$. 
Then, $\sigma(X) = \sigma(X_1, \dots, X_n)$.   
\end{block}
\end{frame}
\begin{frame}{Event space generated by random vectors}
\begin{Proof}
\begin{itemize}
\item <1-> $\sigma(X)$ is generated by the family $(A_X(x): x \in \R^n)$ and $\sigma(X_1, \dots, X_n)$ is generated by the family $(A_{X_i}(x_i): x_i \in \R, i \in [n])$. 
\item <2-> $A_{X_i}(x_i) = A_X(x)$ for $x = (\infty, \dots, x_i, \dots, \infty)$, and hence $\sigma(X_1, \dots, X_n) \subseteq \sigma(X)$.
\item <3-> $A_X(x) = \cap_{i=1}^nA_{X_i}(x_i)$, and hence $\sigma(X) \subseteq \sigma(X_1, \dots, X_n)$. 
\end{itemize}
\end{Proof}

\begin{Example}[Tuple of indicators]<4->
Consider a probability space $(\Omega, \sF, P)$, a finite $n \in \N$, and events $A_1, \dots, A_n \in \sF$, 
that define the random vector $X \triangleq (\Ind{A_1}, \dots, \Ind{A_n})$. 
The $\sigma(X) = \sigma(A_i: i \in [n])$. 
\end{Example}
\end{frame}

\begin{frame}{Independence of random variables}
\begin{Definition}
A family of collections of events $(\sA_i \subseteq \sF: i \in I)$ is called independent,
if for any finite set $F \subseteq I$ and $A_i \in \sA_i$ for all $i \in F$, we have 
\EQ{
P(\cap_{i \in F}A_i) = \prod_{i \in F}P(A_i).
} 
\end{Definition}
\end{frame}
\begin{frame}{Independence of random variables}
\begin{Definition}[Independent and identically distributed]
\begin{itemize}
\item <1-> A random vector $X: \Omega \to \R^n$ defined on $(\Omega, \sF, P)$ is called \textbf{independent} if 
\EQ{
F_X(x) = \prod_{i=1}^nF_{X_i}(x_i),\text{ for all } x \in \R^n.
}
\item <2-> The random vector $X$ is called \textbf{identically distributed} if each of its components have the identical marginal distribution, i.e.
\EQ{
F_{X_i} = F_{X_1}, \text{ for all } i \in [n]. 
}
\end{itemize}
\end{Definition}
\end{frame}
\begin{frame}{Independence of random vectors}
\begin{block}{Reminder}
\begin{itemize}
\item <1-> Events $(A_{X_i}(x_i): i \in [n])$ are independent for any $x \in \R^n$. 
Define families $\sA_i \triangleq (A_{X_i}(x): x \in \R)$ for all $i \in [n]$. Note $(\sA_1, \dots, \sA_n)$ are mutually independent. 
\item <2-> In general, if two collection of events are mutually independent, then the event space generated by them are independent. 
This can be proved using Dynkin's $\pi$-$\lambda$ Theorem. 
\end{itemize}
\end{block}
\end{frame}
\begin{frame}{Independence of random vectors}
\begin{block}{Theorem}
For an independent random vector $X: \Omega \to \R^n$ defined on a probability space $(\Omega, \sF, P)$, 
the event spaces generated by its components $(\sigma(X_i): i \in [n])$ are independent.   
\end{block}
\end{frame}
\begin{frame}{Independence of random vectors}
\begin{Proof}
\begin{itemize}
\item Define a family of events $\sA_i \triangleq (X_i^{-1}(-\infty, x]: x\in \R)$ for each $i \in [n]$. 
\item The families $(\sA_i \subseteq \sF: i \in [n])$ are mutually independent. 
\item Since $\sigma(\sA_i) = \sigma(X_i)$, 
the result follows from the previous remark. 
\end{itemize}
\end{Proof}
\end{frame}
\begin{frame}{Independence of random vectors}
\begin{Definition}[Independent random vectors]
$X:\Omega \to \R^n$ and $Y: \Omega\to\R^m$ defined on same $(\Omega, \sF, P)$ are independent, 
if the collection of events $(A_X(x): x \in \R^n)$ and $(A_Y(y): y \in \R^m)$ are independent,\\ 
where $A_X(x) \triangleq \cap_{i=1}^nX_i^{-1}(-\infty, x_i]$ and $A_Y(y) \triangleq \cap_{i=1}^mY_i^{-1}(-\infty, y_i]$.  
\end{Definition}
\end{frame}
\begin{frame}{Independence of random vectors}
\begin{Example}[Independent random vectors] 
\begin{itemize}
\item <1-> Consider $\sX = \set{(0,0,1),(1,0,0)}\subseteq \R^3$.\\ 
\item <2-> Consider two independent coin tosses, such that $\Omega = \set{H,T}^2, \sF = \cP(\Omega)$ and $P(\omega) = p^{k_2(\omega)}(1-p)^{2-k_2(\omega)}$, where $k_2(\omega) = \sum_{i=1}^2\SetIn{\omega_i = H}$. \\
\item <3->
$&X = (0,0,1)\SetIn{\omega_1 = H} + (1,0,0)\SetIn{\omega_1 = T},$\\
$&Y = (0,0,1)\SetIn{\omega_2 = H} + (1,0,0)\SetIn{\omega_2 = T}.$
\item <4-> $X,Y: \Omega \to \R^3$ are mutually independent random vectors, 
though $X_1, X_3$ are dependent random variables and so are $Y_1, Y_3$.
\end{itemize}
\end{Example}
\end{frame}

\begin{frame}{Discrete random vectors}
\begin{Definition}[Discrete random vectors]
\begin{itemize}
    \item <1-> If a random vector $X: \Omega \to \sX_1 \times \dots  \times \sX_n    \subseteq \R^n$ 
        takes countable values in $\R^n$, 
        then it is called a \textbf{discrete random vector}.
    \item <2-> That is, the range of random vector $X$ is countable,
        and the random vector is completely specified by the \textbf{probability mass function}
        \EQ{
        P_X(x) = P(\cap_{i=1}^n\set{X_i = x_i}) \text{ for all }x \in \sX_1\times \dots \times \sX_n.
        }
\end{itemize}
\end{Definition}
\begin{block}{Reminder}<3->
For an independent discrete random vector $X: \Omega \to \R^n$, we have $P_X(x) = \prod_{i=1}^nP_{X_i}(x_i)$ for each $x \in \R^n$. 
\end{block}
\end{frame}
\begin{frame}{Discrete random vectors}
\begin{Example}[Multiple coin tosses] 
\begin{itemize}
\item <1-> For a probability space $(\Omega, \sF, P)$, such that $\Omega = \set{H,T}^n, \sF = 2^\Omega, P(\omega) = \frac{1}{2^n}$ for all $\omega \in \Omega$.  
\item <2-> Consider the random vector $X: \Omega \to \R$ such that $X_i(\omega) = \SetIn{\omega_i = H}$ for each $i \in [n]$. 
\item <3-> $X$ is a bijection from the sample space to the set $\set{0,1}^n$. 
In particular, $X$ is a discrete random variable. 
\item <4-> For any $x \in [0,1]^n$, $N(x) = \sum_{i=1}^n\Ind{[0,1)}(x_i)$ and the joint distribution is 
\EQ{
F_X(x) = \begin{cases}
1, & x_i \ge 1 \text{ for all }i \in [n],\\
\frac{1}{2^{N(x)}}, & x_i \in [0, 1] \text{ for all } i \in [n],\\
0, & x_i < 0 \text{ for some } i \in [n].
\end{cases}
}
\end{itemize}
\end{Example}
\end{frame}

\begin{frame}{Discrete random vectors}
\begin{Example}[Continued]
We can derive the marginal distribution for $i$-th component as 
\EQ{
F_{X_i}(x_i) = 
\begin{cases}
1, & x_i \ge 1,\\
\frac{1}{2}, & x_i \in [0, 1),\\
0, & x_i < 0.
\end{cases}
}
Therefore, it follows that $X$ is an \iid random vector. 
\end{Example}
\end{frame}

\begin{frame}{Continuous random vectors}
\begin{Definition}[Joint density function]<1->
For jointly continuous random vector $X: \Omega \to \R^n$ with joint distribution function $F_X: \R^n \to [0,1]$, 
there exists a \textbf{joint density function} $f_X: \R^n \to [0,\infty)$ such that $f_X(x) = \frac{d^n}{dx_1\dots dx_n}F_X(x)$, and 
\EQ{
F_X(x) = \int_{u_1\le x_1}du_1\dots \int_{u_n \le x_n} du_n f_X(u_1, \dots, u_n).
}
\end{Definition}
\begin{block}{Reminder}<2->
For an independent continuous random vector $X: \Omega \to \R^n$, we have $f_X(x) = \prod_{i=1}^nf_{X_i}(x_i)$ for all $x \in \R^n$. 
\end{block}
\end{frame}

\begin{frame}{Continuous random vectors}
\begin{Example}[Gaussian random vectors] 
\begin{itemize}
\item <1-> For $(\Omega, \sF, P)$, 
Gaussian random vector is a continuous random vector $X: \Omega \to \R^n$ defined by its density function
\EQ{
f_X(x) = \frac{1}{\sqrt{(2\pi)^n\det(\Sigma)}}\exp\left(-\frac{1}{2}(x-\mu)^T\Sigma^{-1}(x-\mu)\right)\text{ for all }x \in \R^n,
}
\item <2-> where, $\mu \in \R^n$ is the mean vector and $\Sigma \in \R^{n \times n}$ is the positive definite covariance matrix . 
\item <3-> The components of the Gaussian random vector are Gaussian random variables, 
which are independent when $\Sigma$ is diagonal matrix and are identically distributed when $\Sigma = \sigma^2 I$. 
\end{itemize}
\end{Example}
\end{frame}
\end{document}
